% META: {"id": "TSPE-PRIMEQ-EX-012", "chapitre": "TSPE-PRIMITIVES-EQDIFF", "type_objet": "exercice", "capacites": ["TSPE-PRIMITIVES-EQDIFF-C2"], "capacites_codes": ["C2"], "methodes": ["M1"], "parcours": 2, "competences": ["modeliser", "calculer"], "duree_min": 8, "mode_creation": "ex_nihilo", "sources_inspiration": [], "fichier_tex": "chapitres/TSPE-PRIMITIVES-EQDIFF/exercices/TSPE-PRIMEQ-EX-012.tex", "corrige_tex": "chapitres/TSPE-PRIMITIVES-EQDIFF/corriges/TSPE-PRIMEQ-CO-012.tex", "status": "approved"}
% BEGIN-VERIFY
% from sympy import symbols, exp, diff, Rational
% t = symbols('t')
% T = 20 + 70*exp(-t/10)
% assert diff(T, t) == -Rational(1,10)*(T - 20)
% assert T.subs(t, 0) == 90
% assert round(float(T.subs(t, 10)), 2) == 45.75
% END-VERIFY
\begin{exercice}{TSPE-PRIMEQ-EX-004}{2}{16}
Un objet a une temperature $T(t)$ (en \degres C, $t$ en minutes) qui verifie la loi de refroidissement de Newton : $T'(t) = -0{,}1\,(T(t) - 20)$, ou $20\degres$C est la temperature ambiante. A l'instant $t=0$, l'objet est a $90\degres$C.
\begin{enumerate}
  \item Mettre l'equation sous la forme $T' = aT + b$ en precisant $a$ et $b$.
  \item Determiner la solution particuliere constante, puis la forme generale des solutions.
  \item Determiner $T(t)$ en utilisant la condition initiale $T(0)=90$.
  \item Calculer $T(10)$ (arrondir a $0{,}01$ pres).
\end{enumerate}
\end{exercice}
